\begin{code}[hide]%
\>[0]\AgdaKeyword{module}\AgdaSpace{}%
\AgdaModule{paper.sections.permutations}\AgdaSpace{}%
\AgdaKeyword{where}\<%
\\
%
\\[\AgdaEmptyExtraSkip]%
\>[0]\AgdaKeyword{open}\AgdaSpace{}%
\AgdaKeyword{import}\AgdaSpace{}%
\AgdaModule{Data.Nat}\AgdaSpace{}%
\AgdaKeyword{using}\AgdaSpace{}%
\AgdaSymbol{(}\AgdaDatatype{ℕ}\AgdaSpace{}%
\AgdaSymbol{;}\AgdaSpace{}%
\AgdaInductiveConstructor{suc}\AgdaSymbol{)}\<%
\\
\>[0]\AgdaKeyword{open}\AgdaSpace{}%
\AgdaKeyword{import}\AgdaSpace{}%
\AgdaModule{Data.Product}\AgdaSpace{}%
\AgdaKeyword{using}\AgdaSpace{}%
\AgdaSymbol{(}\AgdaOperator{\AgdaFunction{\AgdaUnderscore{}×\AgdaUnderscore{}}}\AgdaSpace{}%
\AgdaSymbol{;}\AgdaSpace{}%
\AgdaOperator{\AgdaInductiveConstructor{\AgdaUnderscore{},\AgdaUnderscore{}}}\AgdaSpace{}%
\AgdaSymbol{;}\AgdaSpace{}%
\AgdaField{proj₁}\AgdaSpace{}%
\AgdaSymbol{;}\AgdaSpace{}%
\AgdaField{proj₂}\AgdaSymbol{)}\<%
\\
\>[0]\AgdaKeyword{open}\AgdaSpace{}%
\AgdaKeyword{import}\AgdaSpace{}%
\AgdaModule{Data.Unit}\AgdaSpace{}%
\AgdaKeyword{using}\AgdaSpace{}%
\AgdaSymbol{(}\AgdaRecord{⊤}\AgdaSymbol{)}\<%
\\
\>[0]\AgdaKeyword{open}\AgdaSpace{}%
\AgdaKeyword{import}\AgdaSpace{}%
\AgdaModule{Level}\AgdaSpace{}%
\AgdaKeyword{using}\AgdaSpace{}%
\AgdaSymbol{(}\AgdaFunction{0ℓ}\AgdaSymbol{)}\<%
\\
\>[0]\AgdaKeyword{open}\AgdaSpace{}%
\AgdaKeyword{import}\AgdaSpace{}%
\AgdaModule{Relation.Binary}\AgdaSpace{}%
\AgdaKeyword{using}\AgdaSpace{}%
\AgdaSymbol{(}\AgdaFunction{Rel}\AgdaSymbol{)}\<%
\\
%
\\[\AgdaEmptyExtraSkip]%
\>[0]\AgdaKeyword{open}\AgdaSpace{}%
\AgdaKeyword{import}\AgdaSpace{}%
\AgdaModule{Notations}\AgdaSpace{}%
\AgdaKeyword{using}\AgdaSpace{}%
\AgdaSymbol{(}\AgdaInductiveConstructor{₁₊}\AgdaSpace{}%
\AgdaSymbol{;}\AgdaSpace{}%
\AgdaInductiveConstructor{₂₊}\AgdaSpace{}%
\AgdaSymbol{;}\AgdaSpace{}%
\AgdaInductiveConstructor{₃₊}\AgdaSymbol{)}\<%
\\
\>[0]\AgdaKeyword{open}\AgdaSpace{}%
\AgdaKeyword{import}\AgdaSpace{}%
\AgdaModule{Word.Base}\AgdaSpace{}%
\AgdaKeyword{using}\AgdaSpace{}%
\AgdaSymbol{(}\AgdaDatatype{Word}\AgdaSpace{}%
\AgdaSymbol{;}\AgdaSpace{}%
\AgdaFunction{WRel}\AgdaSpace{}%
\AgdaSymbol{;}\AgdaSpace{}%
\AgdaOperator{\AgdaInductiveConstructor{[\AgdaUnderscore{}]ʷ}}\AgdaSpace{}%
\AgdaSymbol{;}\AgdaSpace{}%
\AgdaInductiveConstructor{ε}\AgdaSpace{}%
\AgdaSymbol{;}\AgdaSpace{}%
\AgdaOperator{\AgdaInductiveConstructor{\AgdaUnderscore{}•\AgdaUnderscore{}}}\AgdaSpace{}%
\AgdaSymbol{;}\AgdaSpace{}%
\AgdaFunction{wmap}\AgdaSymbol{)}\<%
\\
\>[0]\AgdaKeyword{open}\AgdaSpace{}%
\AgdaKeyword{import}\AgdaSpace{}%
\AgdaModule{Presentation.Definitions}\AgdaSpace{}%
\AgdaKeyword{using}\AgdaSpace{}%
\AgdaSymbol{(}\AgdaOperator{\AgdaRecord{\AgdaUnderscore{}IsPresentationOf\AgdaUnderscore{}}}\AgdaSymbol{)}\<%
\\
\>[0]\AgdaKeyword{open}\AgdaSpace{}%
\AgdaKeyword{import}\AgdaSpace{}%
\AgdaModule{Examples.Groups.Symmetric.Tight.Semantics}\AgdaSpace{}%
\AgdaKeyword{using}\AgdaSpace{}%
\AgdaSymbol{(}\AgdaFunction{Permutation′-group}\AgdaSymbol{)}\<%
\\
\>[0]\AgdaKeyword{import}\AgdaSpace{}%
\AgdaModule{Examples.Groups.Symmetric.Theorems}\AgdaSpace{}%
\AgdaSymbol{as}\AgdaSpace{}%
\AgdaModule{SymThm}\<%
\\
%
\\[\AgdaEmptyExtraSkip]%
\>[0]\AgdaKeyword{postulate}\<%
\\
\>[0][@{}l@{\AgdaIndent{0}}]%
\>[2]\AgdaPostulate{proof-elided}\AgdaSpace{}%
\AgdaSymbol{:}\AgdaSpace{}%
\AgdaSymbol{∀}\AgdaSpace{}%
\AgdaSymbol{\{}\AgdaBound{ℓ}\AgdaSymbol{\}}\AgdaSpace{}%
\AgdaSymbol{\{}\AgdaBound{A}\AgdaSpace{}%
\AgdaSymbol{:}\AgdaSpace{}%
\AgdaPrimitive{Set}\AgdaSpace{}%
\AgdaBound{ℓ}\AgdaSymbol{\}}\AgdaSpace{}%
\AgdaSymbol{→}\AgdaSpace{}%
\AgdaBound{A}\<%
\\
%
\\[\AgdaEmptyExtraSkip]%
\>[0]\AgdaKeyword{private}\AgdaSpace{}%
\AgdaKeyword{variable}\<%
\\
\>[0][@{}l@{\AgdaIndent{0}}]%
\>[2]\AgdaGeneralizable{n}\AgdaSpace{}%
\AgdaSymbol{:}\AgdaSpace{}%
\AgdaDatatype{ℕ}\<%
\end{code}
\section{A Worked Example: Permutation Circuits}\label{sec:permutations}

We begin with the smallest complete case study in the library: circuits
built from a single two-wire gate, the swap $\sigma$, which present the
symmetric groups.  The walkthrough exhibits the pipeline --- gates,
relations, cosets, tower, semantics, uniqueness, completeness,
presentation --- that every later case study instantiates, concretely
enough that \S\ref{sec:design} can be read as ``the same thing, with the
gates abstracted''.

\subsection{Gates, generators, and circuits}

A gate set is a family $\aid{Gate} : \mathbb{N} \to \mathsf{Set}$ giving,
for each arity, the gates on that many adjacent wires; for permutations
there is one gate, on two wires.  The generic circuit layer
(\S\ref{sec:circuits}) turns any such family into wire-indexed
\emph{generators}: a gate placed at the bottom of the wires, or a
generator shifted up by one wire.

\begin{code}%
\>[0]\AgdaKeyword{data}\AgdaSpace{}%
\AgdaDatatype{Gate}\AgdaSpace{}%
\AgdaSymbol{:}\AgdaSpace{}%
\AgdaDatatype{ℕ}\AgdaSpace{}%
\AgdaSymbol{→}\AgdaSpace{}%
\AgdaPrimitive{Set}\AgdaSpace{}%
\AgdaKeyword{where}\<%
\\
\>[0][@{}l@{\AgdaIndent{0}}]%
\>[2]\AgdaInductiveConstructor{σ-gate}\AgdaSpace{}%
\AgdaSymbol{:}\AgdaSpace{}%
\AgdaDatatype{Gate}\AgdaSpace{}%
\AgdaNumber{2}\<%
\\
\>[0]\AgdaKeyword{data}\AgdaSpace{}%
\AgdaDatatype{Gen}\AgdaSpace{}%
\AgdaSymbol{:}\AgdaSpace{}%
\AgdaDatatype{ℕ}\AgdaSpace{}%
\AgdaSymbol{→}\AgdaSpace{}%
\AgdaPrimitive{Set}\AgdaSpace{}%
\AgdaKeyword{where}\<%
\\
\>[0][@{}l@{\AgdaIndent{0}}]%
\>[2]\AgdaInductiveConstructor{gate₁}\AgdaSpace{}%
\AgdaSymbol{:}\AgdaSpace{}%
\AgdaDatatype{Gate}\AgdaSpace{}%
\AgdaNumber{1}\AgdaSpace{}%
\AgdaSymbol{→}\AgdaSpace{}%
\AgdaDatatype{Gen}\AgdaSpace{}%
\AgdaSymbol{(}\AgdaInductiveConstructor{₁₊}\AgdaSpace{}%
\AgdaGeneralizable{n}\AgdaSymbol{)}%
\>[34]\AgdaComment{--\ a\ 1-ary\ gate\ on\ the\ bottom\ wire}\<%
\\
%
\>[2]\AgdaInductiveConstructor{gate₂}\AgdaSpace{}%
\AgdaSymbol{:}\AgdaSpace{}%
\AgdaDatatype{Gate}\AgdaSpace{}%
\AgdaNumber{2}\AgdaSpace{}%
\AgdaSymbol{→}\AgdaSpace{}%
\AgdaDatatype{Gen}\AgdaSpace{}%
\AgdaSymbol{(}\AgdaInductiveConstructor{₂₊}\AgdaSpace{}%
\AgdaGeneralizable{n}\AgdaSymbol{)}%
\>[34]\AgdaComment{--\ a\ 2-ary\ gate\ on\ the\ bottom\ two\ wires}\<%
\\
%
\>[2]\AgdaOperator{\AgdaInductiveConstructor{\AgdaUnderscore{}↥}}%
\>[7]\AgdaSymbol{:}\AgdaSpace{}%
\AgdaDatatype{Gen}\AgdaSpace{}%
\AgdaGeneralizable{n}\AgdaSpace{}%
\AgdaSymbol{→}\AgdaSpace{}%
\AgdaDatatype{Gen}\AgdaSpace{}%
\AgdaSymbol{(}\AgdaInductiveConstructor{₁₊}\AgdaSpace{}%
\AgdaGeneralizable{n}\AgdaSymbol{)}%
\>[34]\AgdaComment{--\ the\ same\ generator,\ one\ wire\ higher}\<%
\end{code}
\begin{code}[hide]%
\>[0]\AgdaFunction{Circuit}\AgdaSpace{}%
\AgdaSymbol{:}\AgdaSpace{}%
\AgdaDatatype{ℕ}\AgdaSpace{}%
\AgdaSymbol{→}\AgdaSpace{}%
\AgdaPrimitive{Set}\<%
\\
\>[0]\AgdaFunction{Circuit}\AgdaSpace{}%
\AgdaBound{n}\AgdaSpace{}%
\AgdaSymbol{=}\AgdaSpace{}%
\AgdaDatatype{Word}\AgdaSpace{}%
\AgdaSymbol{(}\AgdaDatatype{Gen}\AgdaSpace{}%
\AgdaBound{n}\AgdaSymbol{)}\<%
\\
%
\\[\AgdaEmptyExtraSkip]%
\>[0]\AgdaFunction{CRel}\AgdaSpace{}%
\AgdaSymbol{:}\AgdaSpace{}%
\AgdaDatatype{ℕ}\AgdaSpace{}%
\AgdaSymbol{→}\AgdaSpace{}%
\AgdaPrimitive{Set₁}\<%
\\
\>[0]\AgdaFunction{CRel}\AgdaSpace{}%
\AgdaBound{n}\AgdaSpace{}%
\AgdaSymbol{=}\AgdaSpace{}%
\AgdaFunction{Rel}\AgdaSpace{}%
\AgdaSymbol{(}\AgdaFunction{Circuit}\AgdaSpace{}%
\AgdaBound{n}\AgdaSymbol{)}\AgdaSpace{}%
\AgdaFunction{0ℓ}\<%
\\
%
\\[\AgdaEmptyExtraSkip]%
\>[0]\AgdaOperator{\AgdaFunction{\AgdaUnderscore{}↑}}\AgdaSpace{}%
\AgdaSymbol{:}\AgdaSpace{}%
\AgdaFunction{Circuit}\AgdaSpace{}%
\AgdaGeneralizable{n}\AgdaSpace{}%
\AgdaSymbol{→}\AgdaSpace{}%
\AgdaFunction{Circuit}\AgdaSpace{}%
\AgdaSymbol{(}\AgdaInductiveConstructor{₁₊}\AgdaSpace{}%
\AgdaGeneralizable{n}\AgdaSymbol{)}\<%
\\
\>[0]\AgdaOperator{\AgdaFunction{\AgdaUnderscore{}↑}}\AgdaSpace{}%
\AgdaSymbol{=}\AgdaSpace{}%
\AgdaFunction{wmap}\AgdaSpace{}%
\AgdaOperator{\AgdaInductiveConstructor{\AgdaUnderscore{}↥}}\<%
\\
%
\\[\AgdaEmptyExtraSkip]%
\>[0]\AgdaKeyword{pattern}\AgdaSpace{}%
\AgdaInductiveConstructor{σ-gen}\AgdaSpace{}%
\AgdaSymbol{=}\AgdaSpace{}%
\AgdaInductiveConstructor{gate₂}\AgdaSpace{}%
\AgdaInductiveConstructor{σ-gate}\<%
\\
%
\\[\AgdaEmptyExtraSkip]%
\>[0]\AgdaFunction{σ}\AgdaSpace{}%
\AgdaSymbol{:}\AgdaSpace{}%
\AgdaFunction{Circuit}\AgdaSpace{}%
\AgdaSymbol{(}\AgdaInductiveConstructor{₂₊}\AgdaSpace{}%
\AgdaGeneralizable{n}\AgdaSymbol{)}\<%
\\
\>[0]\AgdaFunction{σ}\AgdaSpace{}%
\AgdaSymbol{=}\AgdaSpace{}%
\AgdaOperator{\AgdaInductiveConstructor{[}}\AgdaSpace{}%
\AgdaInductiveConstructor{σ-gen}\AgdaSpace{}%
\AgdaOperator{\AgdaInductiveConstructor{]ʷ}}\<%
\\
%
\\[\AgdaEmptyExtraSkip]%
\>[0]\AgdaKeyword{infix}\AgdaSpace{}%
\AgdaNumber{4}\AgdaSpace{}%
\AgdaOperator{\AgdaDatatype{\AgdaUnderscore{}SRel,\AgdaUnderscore{}===\AgdaUnderscore{}}}\AgdaSpace{}%
\AgdaOperator{\AgdaDatatype{\AgdaUnderscore{}VRel,\AgdaUnderscore{}===\AgdaUnderscore{}}}\AgdaSpace{}%
\AgdaOperator{\AgdaPostulate{\AgdaUnderscore{}≈\AgdaUnderscore{}}}\<%
\\
\>[0]\AgdaKeyword{infixr}\AgdaSpace{}%
\AgdaNumber{10}\AgdaSpace{}%
\AgdaOperator{\AgdaInductiveConstructor{σ•\AgdaUnderscore{}}}\<%
\\
%
\\[\AgdaEmptyExtraSkip]%
\>[0]\AgdaKeyword{private}\AgdaSpace{}%
\AgdaKeyword{variable}\<%
\\
\>[0][@{}l@{\AgdaIndent{0}}]%
\>[2]\AgdaGeneralizable{w}\AgdaSpace{}%
\AgdaGeneralizable{v}\AgdaSpace{}%
\AgdaSymbol{:}\AgdaSpace{}%
\AgdaFunction{Circuit}\AgdaSpace{}%
\AgdaGeneralizable{n}\<%
\\
%
\\[\AgdaEmptyExtraSkip]%
\>[0]\AgdaKeyword{postulate}\<%
\\
\>[0][@{}l@{\AgdaIndent{0}}]%
\>[2]\AgdaOperator{\AgdaPostulate{\AgdaUnderscore{}≈\AgdaUnderscore{}}}\AgdaSpace{}%
\AgdaSymbol{:}\AgdaSpace{}%
\AgdaFunction{Circuit}\AgdaSpace{}%
\AgdaGeneralizable{n}\AgdaSpace{}%
\AgdaSymbol{→}\AgdaSpace{}%
\AgdaFunction{Circuit}\AgdaSpace{}%
\AgdaGeneralizable{n}\AgdaSpace{}%
\AgdaSymbol{→}\AgdaSpace{}%
\AgdaPrimitive{Set}\<%
\end{code}

Here $\aid{₁₊}\;n$ and $\aid{₂₊}\;n$ are the library's patterns for
$1 + n$ and $2 + n$, and $n$ is an implicit argument, so
$\aid{gate₂}\;\aid{σ-gate}$ is a generator at any width from two upward.
An $n$-wire \emph{circuit} is a word of generators, $\aid{Circuit}\;n =
\aid{Word}\;(\aid{Gen}\;n)$ (\S\ref{sec:words}): a single generator
$[\,g\,]$\aid{ʷ}, the empty circuit \aid{ε} --- $n$ bare wires, the
identity circuit at that width --- or a sequential composition
$c \aidop{•} c'$.  We write $\sigma$ for the one-letter circuit
$[\,\aid{gate₂}\;\aid{σ-gate}\;]$\aid{ʷ} and $c\;{\uparrow}$ for $c$ shifted
up one wire ($\aid{{\_}↑} = \aid{wmap}\;\aid{{\_}↥}$).  Read
diagrammatically, with wires drawn bottom-up and composition left to
right, $\sigma$ is a crossing of the bottom two wires, $\sigma\;{\uparrow}$
the same crossing one wire higher, and $\sigma \aidop{•} \sigma\;{\uparrow}
\aidop{•} \sigma$ the three-wire braid on the left of the
\aid{yang-baxter} equation drawn beside its axiom below.

\subsection{Relations: two axioms, structure for free}

The group-specific relations are an inductive family indexed by the wire
count.  Each axiom is stated once, with its gates at the bottom of the
wires and the width left open (the implicit $n$), so that one constructor
covers every width at which the axiom fits; the structural rule
\aid{cong↑} below moves it to any height.  (In Agda, underscores in a
name mark argument positions: the mixfix \aid{{\_}SRel,{\_}==={\_}} is
written $n\;\aid{SRel,}\;w \aidrel{===} v$, ``$w = v$ is an axiom at $n$
wires''.)

\begin{center}
\begin{minipage}[c]{0.58\linewidth}
\begingroup\setlength{\mathindent}{0pt}
\begin{code}%
\>[0]\AgdaKeyword{data}\AgdaSpace{}%
\AgdaOperator{\AgdaDatatype{\AgdaUnderscore{}SRel,\AgdaUnderscore{}===\AgdaUnderscore{}}}\AgdaSpace{}%
\AgdaSymbol{:}\AgdaSpace{}%
\AgdaSymbol{(}\AgdaBound{n}\AgdaSpace{}%
\AgdaSymbol{:}\AgdaSpace{}%
\AgdaDatatype{ℕ}\AgdaSymbol{)}\AgdaSpace{}%
\AgdaSymbol{→}\AgdaSpace{}%
\AgdaFunction{WRel}\AgdaSpace{}%
\AgdaSymbol{(}\AgdaDatatype{Gen}\AgdaSpace{}%
\AgdaBound{n}\AgdaSymbol{)}\AgdaSpace{}%
\AgdaKeyword{where}\<%
\\
\>[0][@{}l@{\AgdaIndent{0}}]%
\>[2]\AgdaInductiveConstructor{order}%
\>[14]\AgdaSymbol{:}\AgdaSpace{}%
\AgdaSymbol{(}\AgdaInductiveConstructor{₂₊}\AgdaSpace{}%
\AgdaGeneralizable{n}\AgdaSymbol{)}\AgdaSpace{}%
\AgdaOperator{\AgdaDatatype{SRel,}}\AgdaSpace{}%
\AgdaFunction{σ}\AgdaSpace{}%
\AgdaOperator{\AgdaInductiveConstructor{•}}\AgdaSpace{}%
\AgdaFunction{σ}\AgdaSpace{}%
\AgdaOperator{\AgdaDatatype{===}}\AgdaSpace{}%
\AgdaInductiveConstructor{ε}\<%
\\
%
\>[2]\AgdaInductiveConstructor{yang-baxter}\AgdaSpace{}%
\AgdaSymbol{:}\AgdaSpace{}%
\AgdaSymbol{(}\AgdaInductiveConstructor{₃₊}\AgdaSpace{}%
\AgdaGeneralizable{n}\AgdaSymbol{)}\AgdaSpace{}%
\AgdaOperator{\AgdaDatatype{SRel,}}\AgdaSpace{}%
\AgdaFunction{σ}\AgdaSpace{}%
\AgdaOperator{\AgdaInductiveConstructor{•}}\AgdaSpace{}%
\AgdaFunction{σ}\AgdaSpace{}%
\AgdaOperator{\AgdaFunction{↑}}\AgdaSpace{}%
\AgdaOperator{\AgdaInductiveConstructor{•}}\AgdaSpace{}%
\AgdaFunction{σ}\AgdaSpace{}%
\AgdaOperator{\AgdaDatatype{===}}\AgdaSpace{}%
\AgdaFunction{σ}\AgdaSpace{}%
\AgdaOperator{\AgdaFunction{↑}}\AgdaSpace{}%
\AgdaOperator{\AgdaInductiveConstructor{•}}\AgdaSpace{}%
\AgdaFunction{σ}\AgdaSpace{}%
\AgdaOperator{\AgdaInductiveConstructor{•}}\AgdaSpace{}%
\AgdaFunction{σ}\AgdaSpace{}%
\AgdaOperator{\AgdaFunction{↑}}\<%
\end{code}
\endgroup
\end{minipage}\hspace{2.5em}%
\begin{minipage}[c]{0.28\linewidth}\centering
  \ufig[0.35]{sn-order}\\[4pt]
  \ufig[0.35]{sn-yb}
\end{minipage}
\end{center}

\noindent These are the Coxeter presentation of $S_n$ in circuit clothing:
$\sigma^2 = \varepsilon$ and the braid (Yang--Baxter) relation.  What
about the relations every circuit theory needs, whatever its gates?  By a
\emph{circuit theory} we mean a gate set with gate-specific relations,
presenting for each width $n$ a monoid of $n$-wire circuits under
sequential composition.  Such circuits are the fixed-width slices of a
symmetric monoidal category (a PROP), and the monoidal structure leaves
exactly two traces at fixed width: the shift $c \mapsto c\;{\uparrow}$ is
``tensoring with an identity wire'', so equations must be preserved by
it, and the interchange law says that gates on disjoint wires commute.
(The version of the library developed since this paper was written also
admits $0$-ary gates --- global scalars, available at every width --- and
adds the rule they need: a scalar commutes with every generator and is
identified with its own shift.)  These are the structural rules every
circuit theory shares --- what a
pen-and-paper presentation leaves implicit in ``a congruence compatible
with composition and tensor product'' --- and the circuit layer adds them
generically: \aid{Lift-Relation} (\S\ref{sec:circuits}) extends any
gate-specific family to the full relation set \aid{{\_}VRel,{\_}==={\_}}:

\begin{code}%
\>[0]\AgdaKeyword{data}\AgdaSpace{}%
\AgdaOperator{\AgdaDatatype{\AgdaUnderscore{}VRel,\AgdaUnderscore{}===\AgdaUnderscore{}}}\AgdaSpace{}%
\AgdaSymbol{:}\AgdaSpace{}%
\AgdaSymbol{(}\AgdaBound{n}\AgdaSpace{}%
\AgdaSymbol{:}\AgdaSpace{}%
\AgdaDatatype{ℕ}\AgdaSymbol{)}\AgdaSpace{}%
\AgdaSymbol{→}\AgdaSpace{}%
\AgdaFunction{CRel}\AgdaSpace{}%
\AgdaBound{n}\AgdaSpace{}%
\AgdaKeyword{where}\<%
\\
\>[0][@{}l@{\AgdaIndent{0}}]%
\>[2]\AgdaInductiveConstructor{srel}%
\>[8]\AgdaSymbol{:}\AgdaSpace{}%
\AgdaGeneralizable{n}\AgdaSpace{}%
\AgdaOperator{\AgdaDatatype{SRel,}}\AgdaSpace{}%
\AgdaGeneralizable{w}\AgdaSpace{}%
\AgdaOperator{\AgdaDatatype{===}}\AgdaSpace{}%
\AgdaGeneralizable{v}\AgdaSpace{}%
\AgdaSymbol{→}\AgdaSpace{}%
\AgdaGeneralizable{n}\AgdaSpace{}%
\AgdaOperator{\AgdaDatatype{VRel,}}\AgdaSpace{}%
\AgdaGeneralizable{w}\AgdaSpace{}%
\AgdaOperator{\AgdaDatatype{===}}\AgdaSpace{}%
\AgdaGeneralizable{v}\<%
\\
%
\>[2]\AgdaInductiveConstructor{cong↑}\AgdaSpace{}%
\AgdaSymbol{:}\AgdaSpace{}%
\AgdaGeneralizable{n}\AgdaSpace{}%
\AgdaOperator{\AgdaDatatype{VRel,}}\AgdaSpace{}%
\AgdaGeneralizable{w}\AgdaSpace{}%
\AgdaOperator{\AgdaDatatype{===}}\AgdaSpace{}%
\AgdaGeneralizable{v}\AgdaSpace{}%
\AgdaSymbol{→}\AgdaSpace{}%
\AgdaSymbol{(}\AgdaInductiveConstructor{₁₊}\AgdaSpace{}%
\AgdaGeneralizable{n}\AgdaSymbol{)}\AgdaSpace{}%
\AgdaOperator{\AgdaDatatype{VRel,}}\AgdaSpace{}%
\AgdaGeneralizable{w}\AgdaSpace{}%
\AgdaOperator{\AgdaFunction{↑}}\AgdaSpace{}%
\AgdaOperator{\AgdaDatatype{===}}\AgdaSpace{}%
\AgdaGeneralizable{v}\AgdaSpace{}%
\AgdaOperator{\AgdaFunction{↑}}\<%
\\
%
\>[2]\AgdaInductiveConstructor{comm₁}\AgdaSpace{}%
\AgdaSymbol{:}\AgdaSpace{}%
\AgdaSymbol{(}\AgdaBound{h}\AgdaSpace{}%
\AgdaSymbol{:}\AgdaSpace{}%
\AgdaDatatype{Gate}\AgdaSpace{}%
\AgdaNumber{1}\AgdaSymbol{)}\AgdaSpace{}%
\AgdaSymbol{(}\AgdaBound{g}\AgdaSpace{}%
\AgdaSymbol{:}\AgdaSpace{}%
\AgdaDatatype{Gen}\AgdaSpace{}%
\AgdaGeneralizable{n}\AgdaSymbol{)}\AgdaSpace{}%
\AgdaSymbol{→}\AgdaSpace{}%
\AgdaSymbol{(}\AgdaInductiveConstructor{₁₊}\AgdaSpace{}%
\AgdaGeneralizable{n}\AgdaSymbol{)}\AgdaSpace{}%
\AgdaOperator{\AgdaDatatype{VRel,}}\AgdaSpace{}%
\AgdaOperator{\AgdaInductiveConstructor{[}}\AgdaSpace{}%
\AgdaBound{g}\AgdaSpace{}%
\AgdaOperator{\AgdaInductiveConstructor{↥}}\AgdaSpace{}%
\AgdaOperator{\AgdaInductiveConstructor{]ʷ}}\AgdaSpace{}%
\AgdaOperator{\AgdaInductiveConstructor{•}}\AgdaSpace{}%
\AgdaOperator{\AgdaInductiveConstructor{[}}\AgdaSpace{}%
\AgdaInductiveConstructor{gate₁}\AgdaSpace{}%
\AgdaBound{h}\AgdaSpace{}%
\AgdaOperator{\AgdaInductiveConstructor{]ʷ}}\AgdaSpace{}%
\AgdaOperator{\AgdaDatatype{===}}\AgdaSpace{}%
\AgdaOperator{\AgdaInductiveConstructor{[}}\AgdaSpace{}%
\AgdaInductiveConstructor{gate₁}\AgdaSpace{}%
\AgdaBound{h}\AgdaSpace{}%
\AgdaOperator{\AgdaInductiveConstructor{]ʷ}}\AgdaSpace{}%
\AgdaOperator{\AgdaInductiveConstructor{•}}\AgdaSpace{}%
\AgdaOperator{\AgdaInductiveConstructor{[}}\AgdaSpace{}%
\AgdaBound{g}\AgdaSpace{}%
\AgdaOperator{\AgdaInductiveConstructor{↥}}\AgdaSpace{}%
\AgdaOperator{\AgdaInductiveConstructor{]ʷ}}\<%
\\
%
\>[2]\AgdaInductiveConstructor{comm₂}\AgdaSpace{}%
\AgdaSymbol{:}\AgdaSpace{}%
\AgdaSymbol{(}\AgdaBound{h}\AgdaSpace{}%
\AgdaSymbol{:}\AgdaSpace{}%
\AgdaDatatype{Gate}\AgdaSpace{}%
\AgdaNumber{2}\AgdaSymbol{)}\AgdaSpace{}%
\AgdaSymbol{(}\AgdaBound{g}\AgdaSpace{}%
\AgdaSymbol{:}\AgdaSpace{}%
\AgdaDatatype{Gen}\AgdaSpace{}%
\AgdaGeneralizable{n}\AgdaSymbol{)}\AgdaSpace{}%
\AgdaSymbol{→}\AgdaSpace{}%
\AgdaSymbol{(}\AgdaInductiveConstructor{₂₊}\AgdaSpace{}%
\AgdaGeneralizable{n}\AgdaSymbol{)}\AgdaSpace{}%
\AgdaOperator{\AgdaDatatype{VRel,}}\AgdaSpace{}%
\AgdaOperator{\AgdaInductiveConstructor{[}}\AgdaSpace{}%
\AgdaBound{g}\AgdaSpace{}%
\AgdaOperator{\AgdaInductiveConstructor{↥}}\AgdaSpace{}%
\AgdaOperator{\AgdaInductiveConstructor{↥}}\AgdaSpace{}%
\AgdaOperator{\AgdaInductiveConstructor{]ʷ}}\AgdaSpace{}%
\AgdaOperator{\AgdaInductiveConstructor{•}}\AgdaSpace{}%
\AgdaOperator{\AgdaInductiveConstructor{[}}\AgdaSpace{}%
\AgdaInductiveConstructor{gate₂}\AgdaSpace{}%
\AgdaBound{h}\AgdaSpace{}%
\AgdaOperator{\AgdaInductiveConstructor{]ʷ}}\AgdaSpace{}%
\AgdaOperator{\AgdaDatatype{===}}\AgdaSpace{}%
\AgdaOperator{\AgdaInductiveConstructor{[}}\AgdaSpace{}%
\AgdaInductiveConstructor{gate₂}\AgdaSpace{}%
\AgdaBound{h}\AgdaSpace{}%
\AgdaOperator{\AgdaInductiveConstructor{]ʷ}}\AgdaSpace{}%
\AgdaOperator{\AgdaInductiveConstructor{•}}\AgdaSpace{}%
\AgdaOperator{\AgdaInductiveConstructor{[}}\AgdaSpace{}%
\AgdaBound{g}\AgdaSpace{}%
\AgdaOperator{\AgdaInductiveConstructor{↥}}\AgdaSpace{}%
\AgdaOperator{\AgdaInductiveConstructor{↥}}\AgdaSpace{}%
\AgdaOperator{\AgdaInductiveConstructor{]ʷ}}\<%
\end{code}

From here on,
$\approx$ denotes the monoid congruence generated by
$n\;\aid{VRel,{\_}==={\_}}$ (\S\ref{sec:presentations}): the least relation
containing the axioms and closed under reflexivity, symmetry, transitivity,
congruence of \aidop{•}, associativity, and the unit laws; deciding it is
the word problem.  Because $\sigma$ is its own inverse and shifts of
invertible generators are invertible, the presented monoid is a group; we
witness this by a \aid{Grouplike} proof (every generator has a left
inverse), by induction on generators.

\subsection{Cosets and the staircase normal form}\label{sec:sn-cosets}

The engine behind the normal form is coset enumeration along the chain
$S_1 < S_2 < \dots < S_n < S_{n+1}$, where $S_n$ sits inside $S_{n+1}$ as
the permutations fixing the bottom wire --- syntactically, the circuits
of the form $c\;{\uparrow}$.  Read as permutations, two circuits lie in the
same right coset of $S_n$ exactly when they bring the same wire to the
bottom position; there are $n{+}1$ choices, so $S_n$ has index $n{+}1$,
and a natural representative of each coset is the \emph{descending
staircase} of swaps that carries that wire down.  In Agda a coset is a
\emph{descriptor}, the length of the staircase as a unary numeral bounded
by the width, and the \emph{section} $[\,\_\,]^{\mathsf c}$ realizes it
as the staircase circuit; the representatives on three, two, and one
wires are drawn on the right:

\begin{center}
\begin{minipage}[c]{0.30\linewidth}
\begingroup\setlength{\mathindent}{0pt}
\begin{code}%
\>[0]\AgdaKeyword{data}\AgdaSpace{}%
\AgdaDatatype{C}\AgdaSpace{}%
\AgdaSymbol{:}\AgdaSpace{}%
\AgdaDatatype{ℕ}\AgdaSpace{}%
\AgdaSymbol{→}\AgdaSpace{}%
\AgdaPrimitive{Set}\AgdaSpace{}%
\AgdaKeyword{where}\<%
\\
\>[0][@{}l@{\AgdaIndent{0}}]%
\>[2]\AgdaInductiveConstructor{ε}%
\>[6]\AgdaSymbol{:}\AgdaSpace{}%
\AgdaDatatype{C}\AgdaSpace{}%
\AgdaGeneralizable{n}\<%
\\
%
\>[2]\AgdaOperator{\AgdaInductiveConstructor{σ•\AgdaUnderscore{}}}\AgdaSpace{}%
\AgdaSymbol{:}\AgdaSpace{}%
\AgdaDatatype{C}\AgdaSpace{}%
\AgdaGeneralizable{n}\AgdaSpace{}%
\AgdaSymbol{→}\AgdaSpace{}%
\AgdaDatatype{C}\AgdaSpace{}%
\AgdaSymbol{(}\AgdaInductiveConstructor{₁₊}\AgdaSpace{}%
\AgdaGeneralizable{n}\AgdaSymbol{)}\<%
\\
\>[0]\AgdaOperator{\AgdaFunction{[\AgdaUnderscore{}]ᶜ}}\AgdaSpace{}%
\AgdaSymbol{:}\AgdaSpace{}%
\AgdaDatatype{C}\AgdaSpace{}%
\AgdaGeneralizable{n}\AgdaSpace{}%
\AgdaSymbol{→}\AgdaSpace{}%
\AgdaFunction{Circuit}\AgdaSpace{}%
\AgdaSymbol{(}\AgdaInductiveConstructor{₁₊}\AgdaSpace{}%
\AgdaGeneralizable{n}\AgdaSymbol{)}\<%
\\
\>[0]\AgdaOperator{\AgdaFunction{[}}\AgdaSpace{}%
\AgdaInductiveConstructor{ε}\AgdaSpace{}%
\AgdaOperator{\AgdaFunction{]ᶜ}}%
\>[11]\AgdaSymbol{=}\AgdaSpace{}%
\AgdaInductiveConstructor{ε}\<%
\\
\>[0]\AgdaOperator{\AgdaFunction{[}}\AgdaSpace{}%
\AgdaOperator{\AgdaInductiveConstructor{σ•}}\AgdaSpace{}%
\AgdaBound{c}\AgdaSpace{}%
\AgdaOperator{\AgdaFunction{]ᶜ}}%
\>[11]\AgdaSymbol{=}\AgdaSpace{}%
\AgdaFunction{σ}\AgdaSpace{}%
\AgdaOperator{\AgdaInductiveConstructor{•}}\AgdaSpace{}%
\AgdaSymbol{(}\AgdaOperator{\AgdaFunction{[}}\AgdaSpace{}%
\AgdaBound{c}\AgdaSpace{}%
\AgdaOperator{\AgdaFunction{]ᶜ}}\AgdaSpace{}%
\AgdaOperator{\AgdaFunction{↑}}\AgdaSymbol{)}\<%
\end{code}
\endgroup
\end{minipage}\hspace{2.5em}%
\begin{minipage}[c]{0.42\linewidth}\centering
\begin{tabular}{@{}l@{\;\;}ccc@{}}
  $\aid{C}\;2$: & \ufig[0.35]{c2-0} & \ufig[0.35]{c2-1} & \ufig[0.35]{c2-2}\\
  & $\varepsilon$ & $\sigma{\bullet}\:\varepsilon$ &
    $\sigma{\bullet}\:\sigma{\bullet}\:\varepsilon$\\[4pt]
  $\aid{C}\;1$, $\aid{C}\;0$: & \ufig[0.35]{c1-0} & \ufig[0.35]{c1-1} & \ufig[0.35]{c0-0}\\
  & $\varepsilon$ & $\sigma{\bullet}\:\varepsilon$ & $\varepsilon$
\end{tabular}
\end{minipage}
\end{center}

\noindent So $\aid{C}\;n$ has exactly the $n{+}1$ elements drawn above
for $n = 2, 1, 0$, matching the index of $S_n$ in $S_{n+1}$.  Why this yields
a normal form is the classical argument of \S\ref{sec:preliminaries}:
every $f \in S_{n+1}$ lies in one coset, say with representative $c_n$,
so $f = f' \cdot c_n$ with $f' \in S_n$, and recursing on $f'$ writes $f$
uniquely as $c_1 c_2 \cdots c_n$ with $c_k \in \aid{C}\;k$ --- a numeral
in the \emph{factorial number system}, one digit per level, since
$|S_{n+1}| = (n{+}1) \cdot n \cdots 2$.

\emph{Computing} this factorization is the job of the coset table, a
right action that pushes one generator past a staircase and returns a
residual circuit on the wires above together with the updated coset.
Pushing $\sigma$ into a staircase can lengthen it, absorb it ($\sigma^2 =
\varepsilon$), or pass it through and shift it up (Yang--Baxter), or
commute past; for a generator that does not touch the bottom wire we can
recurse.  What is
left behind is nothing, one swap on the wires above, or in the other case
studies a proper circuit on the wires above, which is why the table
returns a circuit on one wire fewer rather than a generator
(Figure~\ref{fig:pushing} draws the four cases):

\begin{center}
\begin{minipage}[c]{0.72\linewidth}
\begingroup\setlength{\mathindent}{0pt}
\begin{code}%
\>[0]\AgdaFunction{ract}\AgdaSpace{}%
\AgdaSymbol{:}\AgdaSpace{}%
\AgdaDatatype{C}\AgdaSpace{}%
\AgdaSymbol{(}\AgdaInductiveConstructor{₁₊}\AgdaSpace{}%
\AgdaGeneralizable{n}\AgdaSymbol{)}\AgdaSpace{}%
\AgdaSymbol{→}\AgdaSpace{}%
\AgdaDatatype{Gen}\AgdaSpace{}%
\AgdaSymbol{(}\AgdaInductiveConstructor{₂₊}\AgdaSpace{}%
\AgdaGeneralizable{n}\AgdaSymbol{)}\AgdaSpace{}%
\AgdaSymbol{→}\AgdaSpace{}%
\AgdaFunction{Circuit}\AgdaSpace{}%
\AgdaSymbol{(}\AgdaInductiveConstructor{₁₊}\AgdaSpace{}%
\AgdaGeneralizable{n}\AgdaSymbol{)}\AgdaSpace{}%
\AgdaOperator{\AgdaFunction{×}}\AgdaSpace{}%
\AgdaDatatype{C}\AgdaSpace{}%
\AgdaSymbol{(}\AgdaInductiveConstructor{₁₊}\AgdaSpace{}%
\AgdaGeneralizable{n}\AgdaSymbol{)}\<%
\\
\>[0]\AgdaFunction{ract}\AgdaSpace{}%
\AgdaInductiveConstructor{ε}\AgdaSpace{}%
\AgdaInductiveConstructor{σ-gen}\AgdaSpace{}%
\AgdaSymbol{=}\AgdaSpace{}%
\AgdaInductiveConstructor{ε}\AgdaSpace{}%
\AgdaOperator{\AgdaInductiveConstructor{,}}\AgdaSpace{}%
\AgdaOperator{\AgdaInductiveConstructor{σ•}}\AgdaSpace{}%
\AgdaInductiveConstructor{ε}%
\>[36]\AgdaComment{--\ grows}\<%
\\
\>[0]\AgdaFunction{ract}\AgdaSpace{}%
\AgdaSymbol{(}\AgdaOperator{\AgdaInductiveConstructor{σ•}}\AgdaSpace{}%
\AgdaInductiveConstructor{ε}\AgdaSymbol{)}\AgdaSpace{}%
\AgdaInductiveConstructor{σ-gen}\AgdaSpace{}%
\AgdaSymbol{=}\AgdaSpace{}%
\AgdaInductiveConstructor{ε}\AgdaSpace{}%
\AgdaOperator{\AgdaInductiveConstructor{,}}\AgdaSpace{}%
\AgdaInductiveConstructor{ε}%
\>[36]\AgdaComment{--\ absorbs\ (σ²\ =\ ε)}\<%
\\
\>[0]\AgdaFunction{ract}\AgdaSpace{}%
\AgdaSymbol{(}\AgdaOperator{\AgdaInductiveConstructor{σ•}}\AgdaSpace{}%
\AgdaOperator{\AgdaInductiveConstructor{σ•}}\AgdaSpace{}%
\AgdaBound{c}\AgdaSymbol{)}\AgdaSpace{}%
\AgdaInductiveConstructor{σ-gen}\AgdaSpace{}%
\AgdaSymbol{=}\AgdaSpace{}%
\AgdaFunction{σ}\AgdaSpace{}%
\AgdaOperator{\AgdaInductiveConstructor{,}}\AgdaSpace{}%
\AgdaOperator{\AgdaInductiveConstructor{σ•}}\AgdaSpace{}%
\AgdaOperator{\AgdaInductiveConstructor{σ•}}\AgdaSpace{}%
\AgdaBound{c}%
\>[36]\AgdaComment{--\ passes\ (Yang-Baxter)}\<%
\\
\>[0]\AgdaFunction{ract}\AgdaSpace{}%
\AgdaInductiveConstructor{ε}\AgdaSpace{}%
\AgdaSymbol{(}\AgdaBound{g}\AgdaSpace{}%
\AgdaOperator{\AgdaInductiveConstructor{↥}}\AgdaSymbol{)}\AgdaSpace{}%
\AgdaSymbol{=}\AgdaSpace{}%
\AgdaOperator{\AgdaInductiveConstructor{[}}\AgdaSpace{}%
\AgdaBound{g}\AgdaSpace{}%
\AgdaOperator{\AgdaInductiveConstructor{]ʷ}}\AgdaSpace{}%
\AgdaOperator{\AgdaInductiveConstructor{,}}\AgdaSpace{}%
\AgdaInductiveConstructor{ε}%
\>[36]\AgdaComment{--\ commutes\ past\ /\ recurses}\<%
\\
\>[0]\AgdaFunction{ract}\AgdaSpace{}%
\AgdaSymbol{(}\AgdaOperator{\AgdaInductiveConstructor{σ•}}\AgdaSpace{}%
\AgdaBound{c}\AgdaSymbol{)}\AgdaSpace{}%
\AgdaSymbol{(}\AgdaBound{g}\AgdaSpace{}%
\AgdaOperator{\AgdaInductiveConstructor{↥}}\AgdaSymbol{)}\AgdaSpace{}%
\AgdaSymbol{=}\AgdaSpace{}%
\AgdaField{proj₁}\AgdaSpace{}%
\AgdaSymbol{(}\AgdaFunction{ract}\AgdaSpace{}%
\AgdaBound{c}\AgdaSpace{}%
\AgdaBound{g}\AgdaSymbol{)}\AgdaSpace{}%
\AgdaOperator{\AgdaFunction{↑}}\AgdaSpace{}%
\AgdaOperator{\AgdaInductiveConstructor{,}}\AgdaSpace{}%
\AgdaOperator{\AgdaInductiveConstructor{σ•}}\AgdaSpace{}%
\AgdaSymbol{(}\AgdaField{proj₂}\AgdaSpace{}%
\AgdaSymbol{(}\AgdaFunction{ract}\AgdaSpace{}%
\AgdaBound{c}\AgdaSpace{}%
\AgdaBound{g}\AgdaSymbol{))}%
\>[64]\AgdaComment{--\ recurses}\<%
\end{code}
\endgroup
\end{minipage}\hspace{1em}%
\begin{minipage}[c]{0.24\linewidth}\centering
  \ufig[0.35]{ract-l}\;\scalebox{0.7}{$\approx$}\;\ufig[0.35]{ract-r}
\end{minipage}
\end{center}

\begin{figure}[h]
\centering
\begin{tabular}{@{}r@{\,}c@{\,}l@{\;\;}l@{\qquad}r@{\,}c@{\,}l@{\;\;}l@{}}
  \ufig[0.35]{push4-l} & \scalebox{0.7}{$\rightarrow$} & \ufig[0.35]{push4-r} &
    \small grows &
  \ufig[0.35]{push3-l} & \scalebox{0.7}{$\rightarrow$} & \ufig[0.35]{push3-r} &
    \small absorbs ($\sigma^2 = \varepsilon$)\\[4pt]
  \ufig[0.35]{push2-l} & \scalebox{0.7}{$\rightarrow$} & \ufig[0.35]{push2-r} &
    \small passes (Yang--Baxter) &
  \ufig[0.35]{push1-l} & \scalebox{0.7}{$\rightarrow$} & \ufig[0.35]{push1-r} &
    \small far-commutation
\end{tabular}
\caption{Pushing a swap through a staircase: the four cases of the coset
table \aid{ract} as circuit rewrites, in the order of its clauses.}
\label{fig:pushing}
\end{figure}

\medskip\noindent
\aid{σ-gen} abbreviates $\aid{gate₂}\;\aid{σ-gate}$.  The last clause is the recursion: a
shifted generator meets a staircase by meeting its tail one wire down.
Each clause is certified by the \emph{soundness law}, pictured on the
right and proved once per table by a short equational derivation from
the axioms:

\begin{code}%
\>[0]\AgdaFunction{ract-sound}\AgdaSpace{}%
\AgdaSymbol{:}\AgdaSpace{}%
\AgdaSymbol{∀}\AgdaSpace{}%
\AgdaBound{c}\AgdaSpace{}%
\AgdaBound{b}\AgdaSpace{}%
\AgdaSymbol{→}\AgdaSpace{}%
\AgdaKeyword{let}\AgdaSpace{}%
\AgdaSymbol{(}\AgdaBound{b'}\AgdaSpace{}%
\AgdaOperator{\AgdaInductiveConstructor{,}}\AgdaSpace{}%
\AgdaBound{c'}\AgdaSymbol{)}\AgdaSpace{}%
\AgdaSymbol{=}\AgdaSpace{}%
\AgdaFunction{ract}\AgdaSpace{}%
\AgdaBound{c}\AgdaSpace{}%
\AgdaBound{b}\AgdaSpace{}%
\AgdaKeyword{in}\AgdaSpace{}%
\AgdaOperator{\AgdaFunction{[}}\AgdaSpace{}%
\AgdaBound{c}\AgdaSpace{}%
\AgdaOperator{\AgdaFunction{]ᶜ}}\AgdaSpace{}%
\AgdaOperator{\AgdaInductiveConstructor{•}}\AgdaSpace{}%
\AgdaOperator{\AgdaInductiveConstructor{[}}\AgdaSpace{}%
\AgdaBound{b}\AgdaSpace{}%
\AgdaOperator{\AgdaInductiveConstructor{]ʷ}}\AgdaSpace{}%
\AgdaOperator{\AgdaPostulate{≈}}\AgdaSpace{}%
\AgdaBound{b'}\AgdaSpace{}%
\AgdaOperator{\AgdaFunction{↑}}\AgdaSpace{}%
\AgdaOperator{\AgdaInductiveConstructor{•}}\AgdaSpace{}%
\AgdaOperator{\AgdaFunction{[}}\AgdaSpace{}%
\AgdaBound{c'}\AgdaSpace{}%
\AgdaOperator{\AgdaFunction{]ᶜ}}\<%
\end{code}
\begin{code}[hide]%
\>[0]\AgdaFunction{ract-sound}\AgdaSpace{}%
\AgdaSymbol{=}\AgdaSpace{}%
\AgdaPostulate{proof-elided}\<%
\end{code}

The stateful word traversal \aid{{\_}ᵗ} (\S\ref{sec:words}) extends
\aid{ract} to whole circuits, threading the coset from left to right and
concatenating the residuals; run from the identity coset, it maps an
$(n{+}2)$-wire circuit to an $(n{+}1)$-wire circuit and a coset --- one
level of normalization, the normal-form transport of
\S\ref{sec:engine}.  \emph{Iterating} this is the coset tower, whose
carrier grows by one digit per level:

\begin{center}
\begin{minipage}[c]{0.56\linewidth}
\begingroup\setlength{\mathindent}{0pt}
\begin{code}%
\>[0]\AgdaFunction{NF}\AgdaSpace{}%
\AgdaSymbol{:}\AgdaSpace{}%
\AgdaDatatype{ℕ}\AgdaSpace{}%
\AgdaSymbol{→}\AgdaSpace{}%
\AgdaPrimitive{Set}\<%
\\
\>[0]\AgdaFunction{NF}\AgdaSpace{}%
\AgdaNumber{0}%
\>[11]\AgdaSymbol{=}\AgdaSpace{}%
\AgdaRecord{⊤}\<%
\\
\>[0]\AgdaFunction{NF}\AgdaSpace{}%
\AgdaSymbol{(}\AgdaInductiveConstructor{suc}\AgdaSpace{}%
\AgdaBound{k}\AgdaSymbol{)}\AgdaSpace{}%
\AgdaSymbol{=}\AgdaSpace{}%
\AgdaFunction{NF}\AgdaSpace{}%
\AgdaBound{k}\AgdaSpace{}%
\AgdaOperator{\AgdaFunction{×}}\AgdaSpace{}%
\AgdaDatatype{C}\AgdaSpace{}%
\AgdaBound{k}%
\>[26]\AgdaComment{--\ ⊤\ ×\ C\ 0\ ×\ C\ 1\ ×\ ⋯\ ×\ C\ k}\<%
\end{code}
\endgroup
\end{minipage}\hspace{2.5em}%
\begin{minipage}[c]{0.16\linewidth}\centering
  \ufig[0.35]{fig3_staircase}
\end{minipage}
\end{center}

\noindent so $\aid{NF}\;n$ has exactly $n!$ elements.  The normal-form
function $\aid{nf} : \aid{Circuit}\;n \to \aid{NF}\;n$ runs the tables;
its \emph{section} $\aid{inv-nf} : \aid{NF}\;n \to \aid{Circuit}\;n$
rebuilds a circuit from its digits by concatenating the staircases,
suitably shifted, as drawn on the right for $(\mathsf{tt}, c_1, c_2, c_3)
\in \aid{NF}\;4$ --- one staircase per level, each shifted up to the
wires of its level: a ``staircase of staircases''.

\subsection{Two semantics, uniqueness, and completeness}

So far everything was syntax.  The development interprets circuits twice:
by a \emph{loose} semantics $\aid{⟦\_⟧} : \aid{Circuit}\;n \to
(\mathsf{Fin}\;n \to \mathsf{Fin}\;n)$ into endofunctions of
$\mathsf{Fin}\;n$ ($\sigma$ swaps the bottom two elements), a monoid that
is \emph{not} a group; and by a \emph{tight} semantics into
$\aid{Permutation′-group}\;n$, the standard library's permutations of
$\mathsf{Fin}\;n$ bundled as a \aid{Group} by our \texttt{ForStdlib}
layer (\S\ref{sec:forstdlib}).  \emph{Soundness} --- $w \approx v$
implies $\aid{⟦}w\aid{⟧} \doteq \aid{⟦}v\aid{⟧}$ --- is a small case
analysis over the axioms.  The heart of the matter is the \emph{unique
normal form} property (\S\ref{sec:nfproperty}), that normal forms with
equal denotations are equal: $\aid{⟦}\;\aid{inv-nf}\;u\,\aid{⟧} \doteq
\aid{⟦}\;\aid{inv-nf}\;v\,\aid{⟧}$ implies $u \equiv v$, proved by
induction on the tower, since two staircase sections that agree as
functions send the top point to the same place and so peel off equal top
digits.  Soundness and uniqueness are exactly the premises of the generic
lemma \aid{by-normalization}, which concludes \emph{completeness}:
$\aid{⟦}w\aid{⟧} \doteq \aid{⟦}v\aid{⟧}$ implies $w \approx v$.  The index
module states these for every $n$ (\aid{symmetric-soundness},
\aid{symmetric-completeness}, \aid{symmetric-unique-nf}).  Completeness
for the \emph{loose} semantics is the stronger, perhaps surprising fact:
even read as mere endofunctions, circuits that are equal as functions are
provably equal.  For the tight semantics, soundness, completeness, and
surjectivity (every permutation is a product of adjacent transpositions)
assemble into the \emph{presentation record} of \S\ref{sec:presentations}
--- a \aid{Grouplike} witness with a group isomorphism from the presented
group onto the target whose underlying function is the semantics --- the
one statement we quote in Agda, since every presentation theorem in the
paper has its shape:

\begin{code}%
\>[0]\AgdaFunction{symmetric-presentation}\AgdaSpace{}%
\AgdaSymbol{:}\AgdaSpace{}%
\AgdaSymbol{∀}\AgdaSpace{}%
\AgdaBound{n}\AgdaSpace{}%
\AgdaSymbol{→}\AgdaSpace{}%
\AgdaSymbol{(}\AgdaBound{n}\AgdaSpace{}%
\AgdaOperator{\AgdaDatatype{VRel,\AgdaUnderscore{}===\AgdaUnderscore{}}}\AgdaSymbol{)}\AgdaSpace{}%
\AgdaOperator{\AgdaRecord{IsPresentationOf}}\AgdaSpace{}%
\AgdaSymbol{(}\AgdaFunction{Permutation′-group}\AgdaSpace{}%
\AgdaBound{n}\AgdaSymbol{)}\<%
\end{code}
\begin{code}[hide]%
\>[0]\AgdaFunction{symmetric-presentation}\AgdaSpace{}%
\AgdaBound{n}\AgdaSpace{}%
\AgdaSymbol{=}\AgdaSpace{}%
\AgdaFunction{SymThm.Tight.presentation}\AgdaSpace{}%
\AgdaBound{n}\<%
\end{code}

read: the monoid of swap circuits modulo \{\aid{order},
\aid{yang-baxter}\} and the structural rules \emph{is} the group of
permutations of $\mathsf{Fin}\;n$ (\aid{IsPresentationOf} is the mixfix
record name \aid{{\_}IsPresentationOf{\_}} used infix).  Everything later
in the paper is this pipeline again, with richer gates, longer towers,
and amalgamations where the chain branches.
